\label{sec:results}

\subsection{Experimental Design}

Quantitative results for M.2 land use weights are deferred to Phase~2,
pending download and zonal statistics computation from the NLCD 2021
GeoTIFF.\footnote{Phase~2 will be triggered after the NLCD 2021 GeoTIFF is
downloaded to the bisect data cache.
The data URL is
\url{https://www.mrlc.gov/data/nlcd-2021-land-cover-conus}.
Estimated download size: $\sim$2\,GB compressed.}
This section describes the experimental protocol that will govern Phase~2
execution.

\paragraph{States.}
Three states are selected to represent distinct land use geography profiles:

\begin{itemize}
  \item \textbf{North Carolina ($k=14$).}
    NC provides coastal water complexity (Outer Banks, Pamlico Sound, Albemarle
    Sound), inland development transition zones (Research Triangle residential-commercial
    edge), and agricultural land in the Coastal Plain.
    The hard-cut rule is expected to produce tangible effects on the Pamlico
    Sound adjacencies where mainland tracts and Outer Banks tracts are
    technically contiguous through narrow water passages.
  \item \textbf{Wisconsin ($k=8$).}
    WI provides a test for diffuse land use patterns.
    Milwaukee's urban core has a defined commercial-residential transition;
    the rest of the state is a mixture of agricultural land, small cities,
    and forest.
    This provides a partial contrast to NC: the hard-cut rule may affect
    Green Bay coastal adjacencies; the cosine similarity signal may have
    modest effect in Milwaukee but not statewide.
  \item \textbf{Texas ($k=38$).}
    TX provides the full land use diversity spectrum: Gulf Coast water
    boundaries (Galveston Bay, Matagorda Bay, Corpus Christi Bay), dense
    urban commercial corridors (Houston Loop 610, Dallas Galleria), extensive
    agricultural and ranchland, and desert/scrub in West Texas.
    The hard-cut rule is expected to affect numerous coastal and bay adjacencies.
    The cosine similarity signal may produce strong effects in the Houston and
    Dallas commercial transition zones.
\end{itemize}

All experiments use:
\begin{itemize}
  \item Algorithm: \texttt{standard-bisect} structure
  \item Seed: 42 (single seed; multi-seed variance designated Phase~3)
  \item Blend factor: $\alpha = 0.5$ (\texttt{--land-alpha 0.5})
  \item Year: 2020 decennial census population
  \item Comparison baseline: \texttt{--weights-override geographic}
    on identical plan configuration
\end{itemize}

\subsection{Primary Metric: Land-Use Boundary Coincidence}

The M-track primary metric for M.2 is the \emph{fraction of district
boundaries coinciding with NLCD land-use class transitions}.

A district boundary segment is said to coincide with a land-use transition
if the two tracts separated by that boundary segment have cosine similarity
$\mathrm{sim}(\mathbf{l}_u, \mathbf{l}_v) < 0.4$ under the M.2 vector definition.
Equivalently: the boundary is at a point where the physical land use type
changes substantially (residential-to-commercial, developed-to-natural,
or natural-to-agricultural).

Formally:
\begin{equation}
  \mathrm{LUB} = \frac{\#\{(u,v) \in \partial(\pi) :
    \mathrm{sim}(\mathbf{l}_u, \mathbf{l}_v) < 0.4\}}
  {\#\partial(\pi)}
  \label{eq:lub}
\end{equation}
where $\partial(\pi)$ denotes the set of cut edges in the plan $\pi$.
A higher LUB score means district boundaries fall more consistently at
natural physical transitions, indicating higher physical community coherence.

The prediction for Phase~2:
\begin{itemize}
  \item LUB under land-use weights $\geq$ LUB under geographic weights
    (land-use weights are specifically designed to prefer cuts at transitions).
  \item The improvement in LUB should be at least 20 percentage points on
    NC-14, where the coastal and development-transition geographies are
    well-suited to the signal.
    This target is marked ${}^{\dagger}$Phase~2.
\end{itemize}

\subsection{Secondary Metric: Compactness}

Polsby-Popper compactness \citep{polsby1991third} is the secondary metric.
The prediction is that land-use weights should not reduce compactness:
cutting at physical transition zones tends to produce rounder district shapes
because transition zones are often linear features (shorelines, commercial
strips, ridge lines) that serve as cleaner boundary segments than arbitrary
internal boundaries.

Formally, the M.2 compactness prediction:
\begin{equation}
  \overline{PP}_{\mathrm{landuse}} \geq \overline{PP}_{\mathrm{geographic}}
  \label{eq:pp}
\end{equation}
where $\overline{PP}$ is mean Polsby-Popper across all districts in the plan.
This prediction will be confirmed or refuted in Phase~2.\textsuperscript{$\dagger$}

\subsection{Hard-Cut Validation}

An additional validation metric is specific to the hard-cut rule:

\begin{quote}
  \textit{Zero districts in the final plan cross a majority-water edge
  as defined by equation~\eqref{eq:hardcut}.}
\end{quote}

This is a hard constraint, not a soft metric.
Any plan produced by \texttt{bisect} with \texttt{--weights-override land-use}
must satisfy this constraint by construction (the hard-cut weight of zero
means METIS cannot assign both endpoints of a water edge to the same
partition component under the minimum-cut objective).
Phase~2 will verify this property on the NC Pamlico Sound adjacencies
and the Texas Gulf Coast bay adjacencies as an integration test.

\subsection{Phase 2 Timeline}

All quantitative results in this section are marked ${}^{\dagger}$Phase~2
and will be populated when the following prerequisites are met:

\begin{enumerate}
  \item NLCD 2021 GeoTIFF downloaded to \texttt{data/nlcd/} in the
    bisect data directory.
  \item Zonal statistics computation completed for NC, WI, and TX census
    tracts using the bisect NLCD pipeline (\texttt{bisect fetch --type nlcd}).
  \item Land use Parquet files verified against L0 test invariants
    (Section~\ref{sec:algorithm}).
  \item Single-seed redistricting runs completed under both
    \texttt{land-use} and \texttt{geographic} weight modes.
  \item LUB and Polsby-Popper metrics computed from plan output.
\end{enumerate}
